% Generated draft: requires mathematical and source review.
\documentclass[11pt]{article}
\usepackage[margin=25mm]{geometry}
\usepackage{fontspec}
\setmainfont{Latin Modern Roman}
\usepackage{amsmath,amssymb,mathtools,unicode-math}
\setmathfont{Latin Modern Math}
\usepackage{enumitem,needspace,xurl,hyperref}
\setlength{\parindent}{0pt}
\setlength{\parskip}{5pt}
\setlength{\emergencystretch}{3em}
\newcommand{\N}{\mathbb{N}}
\newcommand{\Z}{\mathbb{Z}}
\newcommand{\Q}{\mathbb{Q}}
\newcommand{\R}{\mathbb{R}}
\newcommand{\C}{\mathbb{C}}
\newcommand{\F}{\mathbb{F}}
\newcommand{\A}{\mathbb{A}}
\newcommand{\PP}{\mathbb{P}}
\newcommand{\E}{\mathbb{E}}
\newcommand{\eps}{\varepsilon}
\newcommand{\dd}{\,\mathrm{d}}
\begin{document}
\section*{UnsolvedMath ID: 30000456}
\textbf{Source problem code:} catalogue entry
\subsection{Definitions and mathematical statement}
In the standard modular version, let $G$ be a finite group, $p$ a prime and $k$ an algebraically closed field of characteristic $p$. A block algebra is $B=kG\,b$ for a primitive central idempotent $b$ of $kG$. Let $D$ be an abelian defect group of this block, and let $B'$ be its Brauer correspondent block in $kN_G(D)$, where
\[
 N_G(D)=\{g\in G:gDg^{-1}=D\}.
\]
Defect groups are the maximal $p$-subgroups detecting the block under the Brauer homomorphism, defined up to conjugacy; they need not be Sylow subgroups of $G$. The Brauer correspondent is the block specified by the usual Brauer first main correspondence for this defect data.

Brou\'e's Abelian Defect Group Conjecture asserts, for every such block,
\[
 D^b(B\text{-}\mathrm{mod})\simeq D^b(B'\text{-}\mathrm{mod}),
\]
a $k$-linear triangulated equivalence between the bounded derived categories of finitely generated modules. This is a derived-equivalence assertion, not merely equality of character counts or a stable equivalence. The source does not demand a splendid, equivariant or otherwise specially structured equivalence; those stronger refinements are not added.

The excerpt does not specify a modular field versus integral modular-system coefficient setting. The display states the customary splitting-field version explicitly; integral lifts require their own coefficient data and cannot be asserted from this sentence alone.

The supplied literature triage checked 21 August 2026 says the source reports several checked cases and an open general conjecture, and that its targeted search located no later general proof. No list of proved block classes or complete later references is supplied. This dated assessment is retained without independent current-status certification.
\subsection{Short English statement}
Must every finite-group block with abelian defect group have the same bounded derived module category as its Brauer correspondent in the defect-group normalizer? The standard modular-field formulation is explicit; no splendid or equivariant refinement is imposed. The supplied assessment reports known cases but no general proof.
\subsection{Sources}
\begin{itemize}
\item Local UnsolvedMath registry, permanent ID 30000456, exact source code catalogue entry; \emph{Brou\'e Abelian Defect-Group Conjecture}. Proposed year 2006, status open; corrected\_verified extraction normalizes only the source label and standard name.
\item Supplied report citation: \emph{Representations of Finite Groups (2007)}, Oberwolfach Reports, DOI 10.4171/owr/2006/15, pp. 895--978, \url{https://doi.org/10.4171/owr/2006/15}. Title year differs from DOI/proposed year.
\item Supplied literature triage checked 2026-08-21, citing only ``original Oberwolfach report'', \url{https://doi.org/10.4171/owr/2006/15}. It reports several checked cases, an open general conjecture, and no later general proof found in the targeted search. The dated OWR-derived curation requests independent verification of time-sensitive status and cited claims.
\end{itemize}
\end{document}
